\documentclass[11pt, a4paper, oneside]{article}
\usepackage[ngerman]{babel}
\usepackage{worksheet}

\begin{document}
	\author{L. Bung}
	\title{Polynomfunktionen: Einführung}
	\subject{Mathematik}
	\class{FHR1}
	\maketitle
	
	\begin{hint}{Polynomfunktionen}
		Addiert man mehrere Potenzfunktionen, so erhält man eine Funktionsgleichung der folgenden Form:
		$$f(x) = a_nx^n + a_{n-1}x^{n-1} + \dots + a_2x^2 + a_1x + a_0$$
		Funktionen dieser Form nennt man \textbf{ganzrationale Funktionen} oder \textbf{Polynomfunktionen}.
		
		Dabei verwendet man die folgenden Begriffe:
		\begin{itemize}
			\item Die höchste Potenz $n$ heißt \emph{Grad der Funktion}.
			\item Die Zahlen $a_n, \dots, a_1, a_0$ heißen \emph{Koeffizienten}.
			\item Der Koeffizient $a_0$ heißt \emph{Absolutglied}.
		\end{itemize}
	\end{hint}
	
	\singletask{Graphen von Polynomfunktionen}
	
	Zeichnen Sie die Graphen der Polynomfunktionen in ein Koordinatensystem.
	
	\begin{multicols}{2}
		\begin{enumerate}[label=\alph*)]
			\item $a(x) = -0,25x^3-0,5x^2+1,25x+1,5$
			\item $b(x) = \frac{1}{5}x^4+\frac{3}{10}x^3-\frac{9}{5}x^2-\frac{17}{10}x+3$
		\end{enumerate}
	\end{multicols}
	
	\checkered[10.5cm]
	
	\partnertask{Verhalten von Polynomfunktionen}
	
	a) Untersuchen Sie, wie sich die Funktionen $f(x) = x^3+x^2+1$ und $g(x) = -3x^5-x^3+3x^2$ verhalten, wenn $x \to +\infty$ oder $x \to -\infty$ geht.
	Was fällt Ihnen auf?
	
	\checkered[2.5cm]
	
	b) Variieren Sie den Funktionsterm des Polynoms $f(x) = x^3+6x^2$ so, dass für $x \to -\infty$ und $x \to +\infty$ gilt, dass $f(x) \to +\infty$.
	
	\checkered[2.5cm]
	
	\begin{hint}{Globaler Verlauf von Polynomfunktionen}
		\phantom{x}\vspace{5.5cm}
	\end{hint}
	
	\begin{goallist}
		...die Graphen von Polynomfunktionen zeichnen. &&&\\
		\hline
		...das Verhalten von Polynomfunktionen für $x \to \infty$ und $x \to -\infty$ beschreiben. &&&\\
		\hline
		...Funktionsterme von Polynomfunktionen ihren Graphen zuordnen. &&&\\
	\end{goallist}
	
\end{document}